% =====================================================================
\chapter{CƠ SỞ LÝ THUYẾT VÀ CÔNG NGHỆ SỬ DỤNG}\label{chap:coso}
% =====================================================================

Chương này trình bày các khái niệm nền tảng của hệ thống nhúng thời gian thực, đặc điểm
của các vi điều khiển, cảm biến và cơ cấu chấp hành được lựa chọn, các giao tiếp ngoại vi,
giao thức mạng và công cụ phát triển. Với mỗi thành phần, báo cáo phân tích lý do lựa chọn
dựa trên yêu cầu của bài toán.

\section{Hệ thống nhúng và hệ điều hành thời gian thực}
Hệ thống nhúng là hệ thống máy tính chuyên dụng được gắn vào một thiết bị lớn hơn để thực
hiện một nhóm chức năng cố định, chạy liên tục với tài nguyên hạn chế và giao tiếp trực tiếp
với thế giới thực qua cảm biến và cơ cấu chấp hành. Thành phần trung tâm là vi điều khiển,
tích hợp CPU, bộ nhớ và các khối ngoại vi như GPIO, bộ định thời, PWM, SPI, ADC và Wi-Fi trên
cùng một chip. Một hệ thống nhúng được gọi là thời gian thực khi tính đúng đắn phụ thuộc cả
vào thời điểm có kết quả. Trong đề tài, thanh chắn phải đóng đúng thời hạn, cảnh báo cháy
phải mở thanh chắn ngay và màn hình phải phản ánh chỗ trống gần như tức thời. Đây là các
ràng buộc thời gian mềm: trễ một chút không gây hỏng hóc nhưng làm giảm chất lượng dịch vụ.

FreeRTOS là nhân hệ điều hành thời gian thực mã nguồn mở được Espressif tích hợp sẵn trong
ESP-IDF. Mỗi công việc được tổ chức thành một tác vụ có ngăn xếp và mức ưu tiên riêng; bộ
lập lịch ưu tiên có chiếm quyền luôn chạy tác vụ sẵn sàng có ưu tiên cao nhất. Tác vụ chờ sự
kiện như gói tin đến hay thông báo từ tác vụ khác không tiêu tốn thời gian CPU. Để chia sẻ
dữ liệu an toàn, FreeRTOS cung cấp vùng găng (critical section) cho các đoạn mã rất ngắn,
mutex cho tài nguyên bị chiếm giữ lâu như bus SPI, và thông báo tác vụ (task notification)
như cơ chế đánh thức nhẹ nhất giữa hai tác vụ. Ngoài ra, bộ hẹn giờ \code{esp\_timer} cho
phép gọi một hàm sau một khoảng thời gian với độ phân giải micro giây.

Máy chủ của hệ thống phải đồng thời nhận gói tin, phục vụ trang web, vẽ màn hình và điều
khiển thanh chắn; nếu viết trong một vòng lặp duy nhất, một thao tác chậm như vẽ toàn màn
hình (khoảng 23\,ms) sẽ làm trễ mọi việc khác. FreeRTOS cho phép việc quan trọng như điều
khiển cổng luôn được thực thi trước việc trang trí như ảnh động. Ngược lại, hai bo thu thập
chỉ thực hiện tuần tự vài thao tác ngắn, nên một vòng lặp không chặn trên nền Arduino là đủ
và đơn giản hơn.

\section{Vi điều khiển}
Hệ thống dùng ba bo mạch với hai loại vi điều khiển. Bảng~\ref{tab:mcu} so sánh ESP32-C3,
được chọn cho máy chủ, với các phương án thường gặp trong đồ án sinh viên.

\begin{table}[H]
\centering\small
\caption{So sánh các phương án vi điều khiển}\label{tab:mcu}
\begin{tabularx}{\linewidth}{@{}lXXXX@{}}
\toprule
Tiêu chí & ESP32-C3 & ESP32 & ESP8266 & Arduino Uno \\ \midrule
CPU & RISC-V 160\,MHz & 2 nhân 240\,MHz & Xtensa 80\,MHz & AVR 16\,MHz \\
SRAM & 400\,KB & 520\,KB & khoảng 80\,KB & 2\,KB \\
Wi-Fi & Có & Có & Có & Không \\
PWM phần cứng & LEDC 14 bit & LEDC & PWM mềm & 8 bit \\
SPI có DMA & Có & Có & Hạn chế & Không \\
RTOS chính thức & FreeRTOS & FreeRTOS & Hạn chế & Không \\
\bottomrule
\end{tabularx}
\end{table}

Phân tích trong Bảng~\ref{tab:mcu} cho thấy yếu tố quyết định là bộ nhớ: riêng ảnh màn hình
$240\times240$ điểm, mỗi điểm 16 bit, đã cần $115\,200$ byte RAM, lớn hơn toàn bộ RAM của
ESP8266 hay Arduino Uno. ESP32-C3 SuperMini có đủ RAM, có DMA cho SPI để CPU không phải chờ
khi truyền ảnh, có khối PWM phần cứng giữ xung servo ổn định khi CPU bận xử lý mạng, và được
hỗ trợ chính thức bởi ESP-IDF cùng FreeRTOS. So với ESP32 bản gốc, ESP32-C3 rẻ và nhỏ hơn
($22\times18$\,mm) trong khi một nhân 160\,MHz vẫn dư năng lực cho bài toán.

Hai bo thu thập dùng ESP8266 vì nhu cầu tính toán rất thấp: đọc vài chân số, một kênh ADC,
giao tiếp SPI với đầu đọc thẻ và gửi gói tin nhỏ. ESP8266 là bo Wi-Fi rẻ, phổ biến và có hệ
sinh thái thư viện Arduino phong phú cho RC522 và DHT11. Việc tách thành hai bo riêng xuất
phát từ giới hạn số chân: màn hình, servo, còi và hai cảm biến cổng đã dùng 10 trong số
khoảng 13 chân khả dụng của ESP32-C3 SuperMini.

\section{Cảm biến và cơ cấu chấp hành}
Các cảm biến và cơ cấu chấp hành được chọn theo ba tiêu chí: giá thành thấp, làm việc ở mức
logic 3,3\,V và tốn ít chân. Bảng~\ref{tab:linhkien} tổng hợp thông số chính và vai trò của
từng linh kiện trong hệ thống.

\begin{table}[H]
\centering\small
\caption{Cảm biến, cơ cấu chấp hành và thông số chính}\label{tab:linhkien}
\begin{tabularx}{\linewidth}{@{}>{\raggedright\arraybackslash}p{3.4cm}>{\raggedright\arraybackslash}p{5.4cm}X@{}}
\toprule
Linh kiện & Thông số chính & Vai trò \\ \midrule
Cảm biến hồng ngoại phản xạ $\times7$ & Tầm 2--30\,cm, bộ so sánh LM393, ngõ ra số tích cực mức thấp & 5 ô đỗ và 2 cổng \\
Cảm biến lửa & Bước sóng 760--1100\,nm, góc khoảng 60°, ngõ ra số có ngưỡng & Phát hiện cháy \\
Cảm biến mực nước & Vùng đo khoảng 40\,mm, ngõ ra tương tự & Phát hiện ngập \\
Đầu đọc RFID MFRC522 $\times2$ & 13,56\,MHz, ISO 14443A, giao tiếp SPI, 3,3\,V & Đọc thẻ ở hai cổng \\
Cảm biến DHT11 & 0--50\,°C, sai số $\pm2$\,°C, giao thức một dây & Nhiệt độ môi trường \\
Servo SG90 $\times2$ & 4,8--6\,V, mô-men khoảng 1,8\,kg$\cdot$cm, điều khiển PWM 50\,Hz & Nâng hạ thanh chắn \\
Màn hình TFT ST7789 & 1,54 inch, $240\times240$, IPS, SPI tối đa khoảng 62\,MHz & Hiển thị tại cổng \\
Còi báo (module) & Tích cực mức thấp & Báo cháy, ngập, xe không thẻ \\
\bottomrule
\end{tabularx}
\end{table}

Theo Bảng~\ref{tab:linhkien}, cảm biến hồng ngoại phản xạ được dùng cho cả ô đỗ và cổng vì
khoảng cách tới thân xe mô hình chỉ vài centimet, nằm trong tầm đo, và mỗi cảm biến chỉ tốn
một chân. Cảm biến siêu âm cho khoảng cách chính xác hơn nhưng tốn hai chân và phải đo thời
gian phản hồi. Cảm biến lửa hồng ngoại phản ứng nhanh hơn cảm biến khói hay nhiệt, phù hợp
với yêu cầu mở lối thoát kịp thời. Đầu đọc MFRC522 đọc được các loại thẻ MIFARE phổ biến như
thẻ sinh viên hay móc khóa từ.

\section{Giao tiếp ngoại vi}
Mỗi thiết bị ngoại vi được nối với vi điều khiển qua một giao tiếp phù hợp với bản chất tín
hiệu của nó. Bảng~\ref{tab:interfaces} liệt kê các giao tiếp được sử dụng.

\begin{table}[H]
\centering\small
\caption{Các giao tiếp ngoại vi được sử dụng}\label{tab:interfaces}
\begin{tabularx}{\linewidth}{@{}p{2.3cm}X>{\raggedright\arraybackslash}p{4.2cm}@{}}
\toprule
Giao tiếp & Nguyên lý & Thiết bị \\ \midrule
GPIO & Chân vào ra số mức 0 hoặc 1, có điện trở kéo lên hoặc kéo xuống nội & 7 cảm biến IR, cảm biến lửa, còi \\
PWM (LEDC) & Khối phần cứng phát xung có chu kỳ cố định và độ rộng xung điều chỉnh được, CPU không phải đảo chân & 2 servo SG90 \\
SPI & Bus nối tiếp đồng bộ gồm SCK, MOSI, MISO và một chân CS cho mỗi thiết bị; có thể dùng DMA & Màn hình ST7789, 2 đầu đọc MFRC522 \\
ADC & Chuyển điện áp tương tự thành số, ESP8266 có một kênh 10 bit & Cảm biến mực nước \\
Một dây & Giao thức riêng của DHT11, dữ liệu 40 bit mã hóa bằng độ rộng xung & Cảm biến nhiệt độ \\
\bottomrule
\end{tabularx}
\end{table}

Điểm chung của các lựa chọn trong Bảng~\ref{tab:interfaces} là ưu tiên khối ngoại vi phần
cứng thay vì điều khiển chân bằng phần mềm. PWM phần cứng giữ xung servo ổn định ngay cả khi
CPU bận vẽ màn hình, còn SPI có DMA truyền 115\,KB dữ liệu ảnh mà CPU không phải sao chép từng
byte. Nếu tạo xung bằng phần mềm, chỉ một ngắt Wi-Fi cũng đủ làm servo rung.

\section{Giao thức mạng và thanh toán}
Các bo mạch và điện thoại của người vận hành liên lạc qua một mạng Wi-Fi chung. Hệ thống sử
dụng các giao thức sau:
\begin{itemize}
  \item Wi-Fi 802.11 b/g/n băng tần 2,4\,GHz ở chế độ trạm, địa chỉ IP do router cấp qua DHCP;
  \item UDP, giao thức không kết nối, độ trễ thấp, cho phép gửi quảng bá (broadcast) tới mọi
        thiết bị trong mạng nên bo thu thập không cần biết địa chỉ của máy chủ;
  \item HTTP cho trang web quản lý và các API dạng JSON, HTTPS khi máy chủ truy vấn SePay;
  \item mDNS để phân giải tên \code{parking.local} trong mạng cục bộ mà không cần máy chủ DNS;
  \item VietQR, chuẩn mã QR chuyển khoản liên ngân hàng của NAPAS dựa trên chuẩn EMVCo, gồm
        chuỗi các trường dạng mã, độ dài và giá trị, kết thúc bằng mã kiểm tra CRC16.
\end{itemize}

UDP được chọn thay cho TCP hay MQTT vì phần lớn dữ liệu là trạng thái được gửi lặp lại mỗi
giây: mất một gói chỉ làm trễ tối đa một giây, trong khi chi phí bắt tay, bộ đệm cho mỗi kết
nối và máy chủ trung gian (broker) là đáng kể trên vi điều khiển. Riêng sự kiện quẹt thẻ
không được gửi lặp lại, nên hệ thống bổ sung cơ chế xác nhận và gửi lại ở tầng ứng dụng
(Mục~\ref{sec:uc02}).

\section{Công cụ phát triển}
Toàn bộ phần mềm sử dụng trong đề tài là miễn phí và mã nguồn mở. Bảng~\ref{tab:software}
liệt kê các thành phần chính.

\begin{table}[H]
\centering\small
\caption{Thành phần phần mềm và vai trò}\label{tab:software}
\begin{tabularx}{\linewidth}{@{}p{4.6cm}X@{}}
\toprule
Thành phần & Vai trò \\ \midrule
ESP-IDF 6.1 và FreeRTOS & Framework cho máy chủ: driver, Wi-Fi, TCP/IP, HTTP, NVS, hệ điều hành \\
\code{esp\_lcd}, LEDC, \code{esp\_timer} & Driver màn hình qua SPI DMA, PWM cho servo, hẹn giờ đóng thanh chắn \\
\code{espressif/mdns}, \code{espressif/qrcode} & Tên \code{parking.local}; mã hóa mã QR thanh toán \\
Arduino core ESP8266, thư viện MFRC522 và DHT & Nền tảng và thư viện thiết bị cho hai bo thu thập \\
HTML, CSS, JavaScript thuần & Trang web quản lý nhúng trong firmware, không phụ thuộc thư viện ngoài \\
Trình biên dịch C/C++ trên máy tính & Chạy bộ kiểm thử tự động không cần phần cứng \\
\bottomrule
\end{tabularx}
\end{table}

Như Bảng~\ref{tab:software} cho thấy, trang web được viết bằng HTML, CSS và JavaScript thuần
và nhúng trực tiếp vào firmware, nên vẫn mở được khi router mất kết nối Internet. Bộ kiểm thử
biên dịch mã nghiệp vụ của máy chủ và mã nguồn hai bo thu thập trên máy tính cùng các thư viện
giả lập phần cứng, giúp kiểm tra logic trong vài giây mà không cần mô hình
(Chương~\ref{chap:kiemthu}).
